\documentclass[11pt]{article}
\usepackage[margin=1in]{geometry}
\usepackage{booktabs}
\usepackage{longtable}
\usepackage{listings}
\usepackage{xcolor}
\usepackage{hyperref}
\usepackage{parskip}

\lstset{
    basicstyle=\ttfamily\small,
    breaklines=true,
    breakatwhitespace=false,
    columns=fullflexible,
    frame=single,
    backgroundcolor=\color{gray!10}
}

\title{Attaching nixnan to a Quantized 7B Medical Dialog Agent:\\
Real Exceptions Found, and a Reproducible Instrumentation Crash}
\author{}
\date{}

\begin{document}
\maketitle

\section{Context}

\texttt{simple\_medical\_agent.py} (in a separate, private project
directory, \texttt{\textasciitilde/repos/newmed} -- not part of this
repository, since it depends on a 14\,GB local model checkpoint and its
own \texttt{venv/}) is a small, standalone medical-dialog chat script
built on \textbf{BioMistral-7B} (\texttt{BioMistral/\allowbreak
BioMistral-7B}), loaded in 4-bit quantization via \texttt{bitsandbytes}
(\texttt{nf4}, double quantization, FP16 compute dtype). It is otherwise
unrelated to the other \texttt{SC26-Tutorial} case studies except that it
uses this repository's own \texttt{nixnan} instrumentation library
(\texttt{external/nixnan/nixnan.so}) to look for floating-point
exceptions inside the model's own CUDA kernels. This report and the logs
behind it live here, in \texttt{biomistral-test/}, since they document a
real \texttt{nixnan} run even though the target program itself does not
live in this repository.

The goal of this exercise: feed the agent a single symptom, \texttt{"dizzy"},
attach \texttt{nixnan} with \texttt{SAMPLING=512}, and let generation run to
completion while collecting exceptions.

\textbf{Headline result: this did not fully succeed.} Every attempt crashed
with a segmentation fault before \texttt{nixnan} printed its own final
aggregate report or the model printed its answer. This document reports
what was found anyway (substantial, real floating-point exception activity)
and treats the crash itself as a finding, not a footnote -- a real,
reproducible robustness limit of NVBit-based instrumentation on a model
this large and architecturally complex, in contrast to every case study in
the sibling \texttt{SC26-Tutorial} project, all of which completed cleanly
under the same tool.

\section{Setup}

\begin{itemize}
  \item \textbf{Model:} \texttt{BioMistral/\allowbreak BioMistral-7B}, 4-bit NF4
    quantization (\texttt{bitsandbytes}), FP16 compute dtype, loaded via
    \texttt{transformers}' \texttt{device\_map="auto"}.
  \item \textbf{Input:} the single symptom \texttt{"dizzy"} (piped via
    stdin, followed by \texttt{quit}/\texttt{n} to exit cleanly), generating
    up to \texttt{max\_new\_tokens=300} with temperature 0.7, top-p 0.9,
    top-k 50.
  \item \textbf{Instrumentation:} \texttt{LD\_PRELOAD} pointed at this
    session's \texttt{nixnan.so} (built from
    \texttt{parfloat/\allowbreak nixnan}, branch \texttt{print\_histo\_\allowbreak instrn} -- the
    same build used throughout \texttt{SC26-Tutorial}), \texttt{SAMPLING=512}
    (samples 1 in every 512 repeat launches of an already-instrumented
    kernel; the \emph{first} launch of every distinct kernel is always
    instrumented, since sampling only gates the runtime enable/disable
    toggle, not the one-time SASS rewrite -- see
    \texttt{SC26-Tutorial/\allowbreak autonixnan-doc.tex}, \S2.2), exception
    checking on by default.
  \item \textbf{Environment:} torch 2.6.0+cu124, transformers 4.57.1,
    bitsandbytes 0.48.1, accelerate 1.10.1, in this directory's own
    \texttt{venv/}.
\end{itemize}

\section{Attempts}

Four attempts were made. Each subsequent one added the previously-observed
crash-adjacent kernel to a \texttt{FUNCTION\_BLACKLIST} file
(\texttt{Logs/function\_blacklist\_final.txt} holds the final, 2-entry
version) to see whether excluding it let the run progress further.

\begin{longtable}{p{1.4cm}p{2.9cm}p{1.3cm}p{2.2cm}p{4.6cm}}
\toprule
\textbf{Att.} & \textbf{Blacklist} & \textbf{Lines} & \textbf{Kernels touched} & \textbf{Last kernel launch logged before the crash} \\
\midrule
\endhead
1 &
  none &
  1581 &
  45 &
  cutlass fused FP16 GEMM+ReLU tensor-core kernel (full name: footnote~\footnotemark[1]) \\
2 &
  none (reproducibility check) &
  1570 &
  45 &
  identical to attempt 1 \\
3 &
  the GEMM+ReLU kernel above &
  1664 &
  91 &
  \texttt{CatArrayBatchedCopy\_\allowbreak aligned16\_\allowbreak contig} (KV-cache concatenation; full name: footnote~\footnotemark[2]) \\
4 &
  both of the above &
  1683 &
  91 &
  same \texttt{CatArrayBatchedCopy} kernel again (see \S\ref{sec:caveat}) \\
\bottomrule
\end{longtable}
\footnotetext[1]{\texttt{cutlass::Kernel2<cutlass\_80\_tensorop\_f16\_s16816gemm\_relu\_f16\_256x64\_32x4\_tn\_align8>}}
\footnotetext[2]{\texttt{CatArrayBatchedCopy\_aligned16\_contig<OpaqueType<4u>, unsigned int, 3, 128, 1>} (template args abbreviated)}

All four crashes were \texttt{SIGSEGV} (``Segmentation fault (core
dumped)''), reported by the shell after the piped-stdin input had already
been fully consumed (the \texttt{printf 'dizzy\textbackslash nquit\textbackslash nn\textbackslash n'} side of the
pipe always finished normally).

Attempts 1 and 2 are the same command run twice, to check whether the
crash was a one-off race or a genuinely reproducible fault: it reproduced
almost exactly (45 kernels touched both times, near-identical exception
counts -- see Table~\ref{tab:exceptions} -- and the identical last-logged
kernel), confirming this is deterministic, not transient.

Each blacklist round did make real forward progress: attempt 3 more than
doubled the number of distinct kernels successfully touched (45
$\rightarrow$ 91) and got well past the point where attempts 1--2 died.
Attempt 4 progressed further still within the log (558 vs.\ 530 distinct
NaN sites, more sampling-related kernels touched -- \texttt{ArgMaxOps},
an exponential-distribution kernel used in nucleus sampling), but still
crashed.

\subsection{A caveat on ``last kernel logged''}
\label{sec:caveat}

\texttt{nixnan} prints \texttt{running kernel [\ldots]} once per distinct
kernel, at launch-interception time, independent of whether that kernel
is actually whitelisted/blacklisted for instrumentation. This means the
last line printed before a crash is not a reliable pointer to the actual
crash site: if the true crash happens during NVBit's one-time SASS
analysis of the \emph{next} distinct kernel, that kernel may crash before
it ever gets to print its own ``running kernel'' line. This is the
likely explanation for why attempt 4's last-logged kernel is the same
\texttt{CatArrayBatchedCopy} kernel as attempt 3's, even though that
specific kernel had already been blacklisted going into attempt 4 --
the log line reflects a normal (blacklisted, uninstrumented, fast) launch
of that kernel, not necessarily the fault. Pinpointing the true crash
site conclusively would require a proper stack trace; this was attempted
(via Ubuntu's \texttt{apport} crash report) but abandoned, since the
resulting core dump was \textbf{57\,GB} -- a 7B-parameter model's
GPU-mapped memory is counted in the crashing process's own address
space -- making it impractical to load into \texttt{gdb} in this
environment.

\section{Exceptions found}
\label{sec:exceptions}

Despite never reaching completion, every attempt captured substantial,
genuine floating-point exception activity, summed across all instrumented
kernels in that run:

\begin{table}[h]
\centering
\begin{tabular}{lrrrrr}
\toprule
\textbf{Attempt} & \textbf{NaN} & \textbf{-Infinity} & \textbf{Subnormal} & \textbf{NaN+Subnormal} & \textbf{Other combos} \\
\midrule
1 & 522 & 214 & 301 & 138 & 1 \\
2 & 518 & 214 & 296 & 137 & 0 \\
3 & 530 & 304 & 229 & 136 & 11 \\
4 & 558 & 304 & 229 & 136 & 2 \\
\bottomrule
\end{tabular}
\caption{Distinct exception sites (instruction + operand combinations),
by category, per attempt. ``Other combos'' groups the rare
\texttt{NaN,Infinity}, \texttt{NaN,-Infinity,subnormal} and
\texttt{NaN,Infinity,subnormal} joint labels. These are counts of
\emph{distinct} sites, not occurrence counts -- each site typically
recurs many times across the model's 32 transformer layers and multiple
attention heads.}
\label{tab:exceptions}
\end{table}

Two distinct exception patterns stand out, both mechanisms already
recognized from the \texttt{SC26-Tutorial} case studies rather than
novel ones:

\begin{itemize}
  \item \textbf{Subnormal-heavy FP16 tensor-core GEMMs.} The bulk of the
    \texttt{subnormal} sites are \texttt{HMMA.16816.F32} tensor-core
    matrix-multiply-accumulate instructions inside
    \texttt{cutlass::Kernel2} template instantiations -- i.e., the
    dequantized 4-bit weights, multiplied through FP16 tensor cores,
    routinely produce subnormal intermediate products. This is consistent
    with 4-bit quantization's narrow dynamic range interacting with FP16's
    already-narrow exponent range.
  \item \textbf{\texttt{-Infinity} from causal-mask softmax.} A large share
    of the \texttt{-infinity} and \texttt{NaN,-infinity} sites are
    \texttt{FADD}/\texttt{FFMA}/\texttt{MUFU.EX2} sequences matching the
    same \texttt{exp(x) = exp2((x - max) \(\times\) log2(e))} Kahan-style
    decomposition already found to interact badly with causal-masked
    \texttt{-inf} values in \texttt{SC26-Tutorial}'s
    \texttt{Karpathy-NGram} case study (see that project's
    \texttt{All-Runs.md}). The mechanism is the same: masked attention
    positions carry a literal \texttt{-inf} into the softmax kernel, and
    intermediate arithmetic on it produces detectable, instruction-level
    \texttt{-Infinity}/\texttt{NaN} values -- resolved to a harmless
    \texttt{0} by the time softmax finishes, in every previously-observed
    case.
\end{itemize}

\section{The crash as a finding}

Every other \texttt{nixnan} case study built in \texttt{SC26-Tutorial} --
spanning custom CUDA kernels, HeCBench ports, a small PyTorch/cuSOLVER
routine, a toy GPT, and a recurrent network -- completed to a clean final
report, even when that required working around a large one-time NVBit
enumeration cost (documented at length for the PyTorch-backed studies).
BioMistral-7B is qualitatively different: it is dramatically larger (7
billion parameters vs.\ a few thousand to a few million in every other
study here), combines 4-bit quantization (\texttt{bitsandbytes}), FlashAttention,
and \texttt{cutlass} fused-epilogue tensor-core kernels in one pipeline, and
reproducibly crashes NVBit's instrumentation machinery partway through a
single forward pass, regardless of \texttt{SAMPLING} (since sampling does
not affect the one-time SASS rewrite that is presumably where the fault
lies).

This is worth stating plainly as a limitation of the tool at its current
state, on workloads at this scale and complexity -- not a failure specific
to this particular medical agent or to the \texttt{"dizzy"} input. The
progressive-blacklisting results (\S3) suggest the fragility is not
confined to one single kernel, but recurs across at least two distinct,
unrelated kernel families (a fused tensor-core GEMM epilogue, and a tensor
concatenation kernel used in KV-cache management), with the true trigger
of the fourth crash not conclusively identified (\S\ref{sec:caveat}).

\section{What was not tried}

Given time constraints, the following mitigations were considered but not
attempted, and would be reasonable next steps:

\begin{itemize}
  \item Continuing the blacklist-and-retry loop further (each round has
    made genuine progress; it is unclear how many more distinct crash
    sites remain).
  \item Running without 4-bit quantization (\texttt{use\_4bit=False}), to
    check whether \texttt{bitsandbytes}' own custom kernels
    (\texttt{kQuantizeBlockwise}/\texttt{kDequantizeBlockwise}) are
    implicated, independent of the \texttt{cutlass}/\texttt{CatArray}
    kernels already isolated.
  \item Capping \texttt{max\_new\_tokens} much lower (e.g.\ 8--16) to
    shrink the number of distinct decode-step kernel launches per run,
    trading completeness of the response for a better chance of reaching
    \texttt{nixnan}'s final report.
  \item Splitting the run \texttt{iter-autonixnan}-style (one kernel
    whitelisted per fresh process) -- impractical for a single coherent
    multi-token \texttt{generate()} call as used here, but could still be
    adapted to at least identify every crash-triggering kernel
    systematically, one process at a time.
  \item Obtaining a proper stack trace via a more targeted core-dump
    strategy (e.g.\ attaching \texttt{gdb} live with a breakpoint, rather
    than post-mortem on a 57\,GB dump).
\end{itemize}

\section{Reproducing this}

This requires \texttt{\textasciitilde/repos/newmed} (the private project
directory holding \texttt{simple\_medical\_agent.py}, its \texttt{venv/},
and the local BioMistral-7B checkpoint) alongside a checkout of this
repository, for \texttt{external/nixnan/nixnan.so}:

\begin{lstlisting}
cd ~/repos/newmed
source venv/bin/activate
NIXNAN_LIB=~/repos/SC26-Tutorial/external/nixnan/nixnan.so
BLACKLIST=~/repos/SC26-Tutorial/biomistral-test/Logs/function_blacklist_final.txt

# Attempt 1 (no blacklist)
printf 'dizzy\nquit\nn\n' | \
  LD_PRELOAD=$NIXNAN_LIB SAMPLING=512 python3 -u simple_medical_agent.py \
  > attempt1-sampling512.log 2>&1

# Attempt 4 (both crash-adjacent kernels blacklisted)
printf 'dizzy\nquit\nn\n' | \
  LD_PRELOAD=$NIXNAN_LIB SAMPLING=512 FUNCTION_BLACKLIST=$BLACKLIST \
  python3 -u simple_medical_agent.py \
  > attempt4-blacklist-relu-gemm-and-catarray.log 2>&1
\end{lstlisting}

All four logs from the original run are in \texttt{Logs/} alongside this
report.

\end{document}
